% Source: https://artofproblemsolving.com/community/c3h3158514_2017_lmt_spring_guts_round_p13p24__lexington_mathematical_tournament
% Stable identity: aops-topic-3158514-q13
A small disk centre $(x,y)$ must be at distance at most $\sqrt2-r $ from each large centre. Thus $ x^2+y^2+2|x|\le(\sqrt2-r)^2-1=:h^2 $. In particular every possible centre lies within radius $ h $ of the origin, so the two centres have separation at most $ 2h $. Disjoint open disks require separation at least $ 2r $, hence $ r\le h $. Squaring yields $ r^2\le1-2\sqrt2r+r^2 $, or $ r\le\dfrac{\sqrt2}{4}$. Equality is achieved with centres $(0,\pm r)$: they are $ 2r $ apart and each has distance $\sqrt{1+r^2}=\sqrt2-r $ from each large centre. Thus $\boxed{r=\dfrac{\sqrt2}{4}}$.
\begin{asy}
usepackage("amsmath");
size(7cm);
real r=sqrt(2)/4;draw(circle((1,0),sqrt(2)));draw(circle((-1,0),sqrt(2)));draw(circle((0,r),r),blue);draw(circle((0,-r),r),blue);dot("$O_1$",(-1,0),W);dot("$O_2$",(1,0),E);dot("$C_1$",(0,r),N);dot("$C_2$",(0,-r),S);
\end{asy}
